\section{Turing Machine --- Build --- Contract --- Foundations}
%%% generated by the blueprint pipeline from TCSlib.Complexity.TuringMachine.Build.Contract.Basic
\subsection*{Turing Machine --- Build --- Contract --- Foundations: Introduction}

Declarations of the Lean module \texttt{TCSlib.Complexity.TuringMachine.Build.Contract.Basic}.

\subsection*{Turing Machine --- Build --- Contract --- Foundations: Results}

\begin{definition}[Machine execution contract]
\lean{Turing.Contract}
\label{Turing.Contract}
\leanok
\uses{Turing.Cfg, Turing.MultiTapeTM, Turing.MultiTapeTM.runFrom}
Let $M$ be a $k$-work-tape Turing machine, let $e$ be an entry state, let $E$ be a set
of possible states (including the halting state), let $I$ be a type of data, and fix an
input string $x$. A contract assigns to each $i\in I$ and configuration $d$ a
precondition, a running time, an endpoint, and integer lower and upper bounds for each
work-tape head. Whenever $d$ starts in state $e$ and satisfies the precondition, running
$M$ for the assigned time reaches the assigned endpoint, whose state lies in $E$; before
that time the run's state is outside $E$, and at every time up to and including the
assigned time each work-tape head lies between its assigned bounds.
\end{definition}

\begin{theorem}[Contract Endpoint Within a Time Bound]
\lean{Turing.Contract.within}
\label{Turing.Contract.within}
\leanok
\uses{Turing.Cfg, Turing.Contract, Turing.MultiTapeTM.runFrom}
\difficulty{1}
Let $M$ be a multi-tape Turing machine and $C$ a contract for $M$ with exit set $E$.
Suppose a starting configuration $d$ is in the entry state and satisfies the
precondition for parameter $i$. For any natural number $B$ at least the prescribed
running time, there is a time $t\leq B$ such that the run's state is outside $E$ at
every earlier time and the configuration at time $t$ is the contract endpoint.
\end{theorem}

\begin{theorem}[Space Bound from a Head Envelope]
\lean{Turing.Contract.spaceUsedByTape_le}
\label{Turing.Contract.spaceUsedByTape_le}
\leanok
\uses{Turing.Cfg, Turing.Contract, Turing.MultiTapeTM.spaceUsedByTape, Turing.MultiTapeTM.spaceUsedByTape_le_card_Icc}
\difficulty{1}
Let $M$ be a machine and $C$ a contract for $M$. For any data value $i$, configuration
$d$, and work-tape index $j$, the number of distinct cells visited by tape $j$ from time
zero through time $C$’s assigned time is at most the natural-number part of $C$’s upper
bound plus one minus its lower bound for that tape; for an integer $z$, this
natural-number part is $\max(z,0)$.
\end{theorem}

\begin{definition}[Positive running time]
\lean{Turing.Contract.Positive}
\label{Turing.Contract.Positive}
\leanok
\uses{Turing.Contract}
A contract has positive running time if, for every data value $i$ and starting
configuration $d$ whose state is the entry state and which satisfies the contract's
precondition, its assigned running time is strictly positive.
\end{definition}

\begin{theorem}[Absolute bound gives an interval]
\lean{Turing.Contract.mem_Icc_of_abs_le}
\label{Turing.Contract.mem_Icc_of_abs_le}
\leanok
\difficulty{3}
For integers $p,q,R$, if $|p-q|\le R$, then $p$ lies in the closed integer interval from
$q-R$ to $q+R$.
\end{theorem}

\begin{definition}[Composition of execution contracts]
\lean{Turing.Contract.trans}
\label{Turing.Contract.trans}
\leanok
\uses{Turing.Contract, Turing.MultiTapeTM, Turing.MultiTapeTM.runFrom, Turing.MultiTapeTM.runFrom_add}
Let $C_1$ be a contract from entry state $e_1$ to the singleton exit state $e_2$, and
let $C_2$ be a contract from $e_2$ to a set $X$ of states. Suppose a map sends each data
value for $C_1$ to a data value for $C_2$, and every endpoint of $C_1$ satisfying its
precondition satisfies the corresponding precondition of $C_2$. Also suppose that before
the end of $C_1$, the run's state is outside $X$. The composed contract uses the
precondition of $C_1$, runs for the sum of the two assigned times, ends at the endpoint
assigned by $C_2$, and bounds each work-tape head between the minimum of the two lower
bounds and the maximum of the two upper bounds.
\end{definition}

\begin{theorem}[Nonhalting before the exit time]
\lean{Turing.Contract.live}
\label{Turing.Contract.live}
\leanok
\uses{Turing.Cfg, Turing.Contract, Turing.MultiTapeTM, Turing.MultiTapeTM.runFrom, Turing.MultiTapeTM.runFrom_state_ne_none_of_le}
\difficulty{4}
Let $C$ be a contract for a machine with entry state $entry$, let $i$ be a datum, and
let $d$ be a starting configuration whose state is $entry$ and which satisfies $C$'s
precondition for $i$. At every natural time $u$ strictly less than $C$'s assigned time
for $(i,d)$, the machine's state is not halted.
\end{theorem}

\begin{definition}[Agreement outside specified cells]
\lean{Turing.Cfg.SameOutside}
\label{Turing.Cfg.SameOutside}
\leanok
\uses{Turing.Cfg}
Let $d$ and $d'$ be configurations with the same input and work tapes indexed by $j<k$,
and for each tape $j$ let $F_j\subseteq\mathbb Z$ be a set of cell positions. The
configurations agree outside these sets if, for every tape $j$ and every position
$q\notin F_j$, the contents of cell $q$ on tape $j$ are the same in $d$ and $d'$.
\end{definition}

\begin{theorem}[Transitivity of tape agreement outside sets]
\lean{Turing.Cfg.SameOutside.trans}
\label{Turing.Cfg.SameOutside.trans}
\leanok
\uses{Turing.Cfg, Turing.Cfg.SameOutside}
\difficulty{3}
Let $d,d',d''$ be configurations with the same number of work tapes, symbol type, and
input, possibly with different state types. Suppose $d$ and $d'$ agree outside position
sets $F_j$, and $d'$ and $d''$ agree outside position sets $F'_j$. Then $d$ and $d''$
agree outside the unions $F_j\cup F'_j$ on each work tape $j$.
\end{theorem}

\begin{theorem}[State mapping preserves tape agreement]
\lean{Turing.Cfg.SameOutside.mapState}
\label{Turing.Cfg.SameOutside.mapState}
\leanok
\uses{Turing.Cfg, Turing.Cfg.SameOutside, Turing.Cfg.mapState}
\difficulty{1}
Let $d$ and $d'$ be configurations with the same number of work tapes, symbol type, and
input, possibly with different state types, and suppose their tapes agree outside
position sets $F_j$. For any function $f$ from the state set of $d'$ to another state
set, the tapes of $d$ and the configuration obtained from $d'$ by applying $f$ to its
state still agree outside those sets.
\end{theorem}

\begin{theorem}[A stored word is preserved outside a footprint]
\lean{Turing.Cfg.SameOutside.keep_word}
\label{Turing.Cfg.SameOutside.keep_word}
\leanok
\uses{Turing.Cfg, Turing.Cfg.SameOutside, Turing.FinTM.bufferTape}
\difficulty{1}
Let $d$ and $d'$ be configurations with the same number of work tapes, symbol type, and
input, possibly with different state types, whose tapes agree outside position sets
$F_j$. Choose a tape $j$, an integer offset $o$, and a finite word $w$ over the symbol
type. Suppose that for every integer $p$ with $-1\le p\le |w|$, the position $o+p$ is
outside $F_j$, and tape $j$ of $d$ at $o+p$ contains the $p$-th symbol of $w$ when $p$
is a valid word index and is blank otherwise. Then tape $j$ of $d'$ has the same
contents at every such position $o+p$.
\end{theorem}

\begin{theorem}[Blank tape preserved outside an empty set]
\lean{Turing.Cfg.SameOutside.keep_blank}
\label{Turing.Cfg.SameOutside.keep_blank}
\leanok
\uses{Turing.Cfg, Turing.Cfg.SameOutside}
\difficulty{1}
Let $d$ and $d'$ be configurations with the same number of work tapes and input,
possibly with different state types, and let $F_j$ be a set of integer positions for
each tape $j$. Suppose the tapes of $d$ and $d'$ agree at every position outside the
corresponding set $F_j$. If $F_j$ is empty for a tape $j$ and that tape of $d$ is blank
at every position, then that tape of $d'$ is blank at every position as well.
\end{theorem}

\begin{definition}[Frame property for a contract]
\lean{Turing.Contract.Frames}
\label{Turing.Contract.Frames}
\leanok
\uses{Turing.Cfg, Turing.Cfg.SameOutside, Turing.Contract}
Let $C$ be a contract and, for each data value $i$, starting configuration $d$, and work
tape $j$, specify a set $F_{i,d,j}$ of cell positions. The contract has the frame
property for these sets if every eligible run ends in a configuration whose work tapes
agree with those of its starting configuration at every position outside the
corresponding set.
\end{definition}

\begin{theorem}[Characterization of frame preservation]
\lean{Turing.Contract.frames_iff}
\label{Turing.Contract.frames_iff}
\leanok
\uses{Turing.Cfg, Turing.Contract, Turing.Contract.Frames}
\difficulty{1}
Let $C$ be a contract and let $F_{i,d,j}\subseteq\mathbb Z$ specify, for each datum $i$,
starting configuration $d$, and tape $j$, a set of cell positions. The endpoint of every
eligible run agrees with its starting configuration on every tape cell outside the
corresponding set if and only if, for every such $i,d,j$ and every $q\notin F_{i,d,j}$,
the endpoint's content at $q$ equals the starting configuration's content there.
\end{theorem}

\begin{definition}[Return-to-state execution contract]
\lean{Turing.ReturnContract}
\label{Turing.ReturnContract}
\leanok
\uses{Turing.Cfg, Turing.MultiTapeTM, Turing.MultiTapeTM.runFrom}
Let $M$ be a $k$-work-tape Turing machine, let $a$ be an anchor state, let $I$ be a type
of data, and fix an input string $x$. A return contract assigns each data value and
configuration a precondition, a strictly positive running time, an endpoint, and integer
lower and upper bounds for each work-tape head. Whenever a configuration starts in state
$a$ and satisfies the precondition, the run reaches its assigned endpoint at the
assigned time, and that endpoint is again in state $a$; at each positive time strictly
before the endpoint, the state is not $a$, and throughout the run each work-tape head
stays within its assigned bounds.
\end{definition}

\begin{definition}[Work-tape head box]
\lean{Turing.HeadBox}
\label{Turing.HeadBox}
\leanok
Given integer-valued lower and upper bounds $\ell_j,h_j$ for each of $k$ work tapes and
a vector of head positions $p_j$, the vector lies in the head box if $\ell_j\le p_j\le
h_j$ for every tape $j$.
\end{definition}

\begin{definition}[Bounded computation segment]
\lean{Turing.Segment}
\label{Turing.Segment}
\leanok
\uses{Turing.Cfg, Turing.MultiTapeTM, Turing.MultiTapeTM.runFrom}
Let $M$ be a $k$-work-tape Turing machine, let $F$ be a set of states, let $P$ be a
property of work-tape head-position vectors, and let $a,b$ be configurations on the same
input. A segment of length $t$ from $a$ to $b$ means that running $M$ from $a$ for $t$
steps reaches $b$, the state is outside $F$ at every time strictly before $t$, and the
head-position vector satisfies $P$ at every time up to and including $t$.
\end{definition}

\begin{theorem}[The zero-step computation segment]
\lean{Turing.Segment.zero}
\label{Turing.Segment.zero}
\leanok
\uses{Turing.Cfg, Turing.Segment}
\difficulty{3}
Let $M$ be a multi-tape Turing machine, let $F$ be a set of states, and let $P$ be a
property of work-tape head positions. If $P$ holds at the work-tape head positions of a
configuration $a$, then the zero-step run from $a$ to itself avoids $F$ at all times
strictly before its endpoint and satisfies $P$ through its endpoint.
\end{theorem}

\begin{theorem}[A one-step computation segment]
\lean{Turing.Segment.one}
\label{Turing.Segment.one}
\leanok
\uses{Turing.Cfg, Turing.MultiTapeTM.step, Turing.Segment}
\difficulty{4}
Let $M$ be a multi-tape Turing machine, let $F$ be a set of states, and let $P$ be a
property of work-tape head positions. If the state of a configuration $a$ is outside
$F$, and $P$ holds for the work-tape head positions of both $a$ and its one-step
successor, then the run from $a$ to that successor in one step avoids $F$ before the
endpoint and satisfies $P$ at both times.
\end{theorem}

\begin{theorem}[Concatenation of computation segments]
\lean{Turing.Segment.trans}
\label{Turing.Segment.trans}
\leanok
\uses{Turing.Cfg, Turing.MultiTapeTM.runFrom_add, Turing.Segment}
\difficulty{5}
Let $M$ be a multi-tape Turing machine, let $F$ be a set of states, and let $P$ be a
property of work-tape head positions. If a run from $a$ reaches $b$ after $t$ steps and
avoids $F$ before time $t$, and a run from $b$ reaches $c$ after $v$ steps and avoids
$F$ before time $v$, with $P$ holding on the work-tape head positions throughout both
segments, then the run from $a$ reaches $c$ after $t+v$ steps, avoids $F$ at every time
before $t+v$, and satisfies $P$ at every time through $t+v$.
\end{theorem}

\begin{theorem}[Monotonicity of segment conditions]
\lean{Turing.Segment.mono}
\label{Turing.Segment.mono}
\leanok
\uses{Turing.Cfg, Turing.Segment}
\difficulty{1}
Let $M$ be a multi-tape Turing machine, let $F$ and $F'$ be sets of states with
$F'\subseteq F$, and let $P$ and $P'$ be properties of work-tape head positions such
that $P(p)$ implies $P'(p)$ for every position tuple $p$. If a run from $a$ reaches $b$
after $t$ steps, avoids $F$ at every time $u<t$, and satisfies $P$ at every time $u\le
t$, then it reaches the same endpoint after the same number of steps, avoids $F'$ at
every time $u<t$, and satisfies $P'$ at every time $u\le t$.
\end{theorem}

\begin{theorem}[Combining segment invariants]
\lean{Turing.Segment.and}
\label{Turing.Segment.and}
\leanok
\uses{Turing.Cfg, Turing.Segment}
\difficulty{1}
Let $M$ be a multi-tape machine, $F$ a set of states, and $P,Q$ predicates on work-tape
head positions. If the same run from $a$ reaches $b$ after $t$ steps, avoids $F$ before
time $t$, satisfies $P$ throughout times $0$ through $t$, and also satisfies $Q$
throughout those times, then it satisfies both $P$ and $Q$ throughout times $0$ through
$t$; thus it is a segment relative to $F$ and the predicate $p\mapsto P(p)\land Q(p)$.
\end{theorem}

\begin{theorem}[Taking an initial segment]
\lean{Turing.Segment.take}
\label{Turing.Segment.take}
\leanok
\uses{Turing.Cfg, Turing.MultiTapeTM.runFrom, Turing.Segment}
\difficulty{3}
Let $M$ be a multi-tape Turing machine, let $F$ be a set of states, and let $P$ be a
property of work-tape head positions. Suppose a run from $a$ reaches $b$ after $t$
steps, avoids $F$ at every time before $t$, and satisfies $P$ at every time through $t$.
For any natural number $u\le t$, the run from $a$ reaches its configuration at time $u$
after $u$ steps, avoids $F$ at every earlier time, and satisfies $P$ at every time
through $u$.
\end{theorem}

\begin{theorem}[Taking a suffix of a segment]
\lean{Turing.Segment.drop}
\label{Turing.Segment.drop}
\leanok
\uses{Turing.Cfg, Turing.MultiTapeTM.runFrom, Turing.MultiTapeTM.runFrom_add, Turing.Segment}
\difficulty{4}
Let $M$ be a multi-tape machine, $F$ a set of states, and $P$ a predicate on work-tape
head positions. If a segment from $a$ to $b$ of length $t$ avoids $F$ before time $t$,
satisfies $P$ at every time from $0$ through $t$, and $u\le t$, then the run starting at
the configuration reached from $a$ after $u$ steps is a segment to $b$ of length $t-u$
relative to $F$ and $P$.
\end{theorem}

\begin{theorem}[Uniqueness of a first hitting time]
\lean{Turing.Segment.time_eq}
\label{Turing.Segment.time_eq}
\leanok
\uses{Turing.Cfg, Turing.Segment}
\difficulty{4}
Let $M$ be a multi-tape Turing machine and let $F$ be a set of states. Suppose a run
from $a$ reaches $b$ after $t$ steps, avoids $F$ at every time before $t$, and the state
of $b$ belongs to $F$. Suppose also that the run from $a$ reaches $b'$ after $t'$ steps,
avoids $F$ at every time before $t'$, and the state of $b'$ belongs to $F$. Then $t=t'$.
\end{theorem}

\begin{theorem}[Prepending one computation step]
\lean{Turing.Segment.cons}
\label{Turing.Segment.cons}
\leanok
\uses{Turing.Cfg, Turing.MultiTapeTM.runFrom_zero, Turing.MultiTapeTM.step, Turing.Segment, Turing.Segment.one, Turing.Segment.trans}
\difficulty{4}
Let $M$ be a multi-tape machine, $F$ a set of states, and $P$ a predicate on work-tape
head positions. Suppose one step of $M$ from $a$ reaches $a_1$, the state of $a$ is
outside $F$, $P$ holds at the work-tape head-position vector of $a$, and there is a
segment from $a_1$ to $b$ of length $t$ relative to $F$ and $P$. Then there is a segment
from $a$ to $b$ of length $1+t$ relative to $F$ and $P$.
\end{theorem}

\begin{theorem}[Extension of a live segment]
\lean{Turing.Segment.live}
\label{Turing.Segment.live}
\leanok
\uses{Turing.Cfg, Turing.MultiTapeTM.runFrom_state_ne_none_of_le, Turing.Segment}
\difficulty{4}
Let $M$ be a multi-tape Turing machine, let $F$ be a set of states (including the
possibility of the halting state), and let $P$ be a property of work-tape head
positions. Suppose a computation from $a$ reaches $b$ after $t$ steps, avoids $F$ at
every time $u<t$, and satisfies $P$ at every time $u\le t$. If $b$ is not halted, then
the same computation segment also avoids $F\cup\{\text{halt}	ext{ state}\}$ at every
time before $t$, while still satisfying $P$ through time $t$.
\end{theorem}

\begin{definition}[Safe computation segment]
\lean{Turing.SafeSegment}
\label{Turing.SafeSegment}
\leanok
\uses{Turing.Cfg, Turing.MultiTapeTM, Turing.Segment}
Let $M$ be a $k$-work-tape Turing machine, let $F$ be a set of states (including
possibly the halting state), and let $P$ be a predicate on the tuple of work-tape head
positions. For configurations $a$ and $b$ on a fixed input and a time $t$, the
computation from $a$ reaches $b$ after exactly $t$ steps, avoids states in $F$ at every
time strictly before $t$, and satisfies $P$ on the work-tape head positions at every
time from $0$ through $t$; moreover, the state of $b$ is not in $F$.
\end{definition}

\begin{theorem}[Characterization of safe computation segments]
\lean{Turing.safeSegment_iff}
\label{Turing.safeSegment_iff}
\leanok
\uses{Turing.Cfg, Turing.MultiTapeTM, Turing.MultiTapeTM.runFrom, Turing.SafeSegment}
\difficulty{4}
Let $M$ be a multi-tape Turing machine, $F$ a set of states (including the possible
halting state), and $P$ a predicate on work-tape head positions. For configurations
$a,b$ and a time $t$, the segment from $a$ to $b$ is safe exactly when $M$ reaches $b$
from $a$ after $t$ steps, every configuration from time $0$ through time $t$ has state
outside $F$, and every such configuration satisfies $P$ of its work-tape head positions.
\end{theorem}

\begin{theorem}[Safe segment of length zero]
\lean{Turing.SafeSegment.zero}
\label{Turing.SafeSegment.zero}
\leanok
\uses{Turing.Cfg, Turing.SafeSegment, Turing.Segment.zero}
\difficulty{1}
Let $M$ be a multi-tape machine, $F$ a set of states, $P$ a predicate on work-tape head
positions, and $a$ a configuration. If the state of $a$ is outside $F$ and $P$ holds at
its work-tape head-position vector, then there is a safe segment from $a$ to itself of
length zero.
\end{theorem}

\begin{theorem}[Concatenation of safe segments]
\lean{Turing.SafeSegment.trans}
\label{Turing.SafeSegment.trans}
\leanok
\uses{Turing.Cfg, Turing.SafeSegment, Turing.Segment.trans}
\difficulty{1}
Let $M$ be a multi-tape machine, $F$ a set of states, and $P$ a predicate on work-tape
head positions. If there is a safe segment from $a$ to $b$ of length $t$ and a safe
segment from $b$ to $c$ of length $v$, both relative to $F$ and $P$, then there is a
safe segment from $a$ to $c$ of length $t+v$.
\end{theorem}

\begin{theorem}[Monotonicity of safe segments]
\lean{Turing.SafeSegment.mono}
\label{Turing.SafeSegment.mono}
\leanok
\uses{Turing.Cfg, Turing.SafeSegment, Turing.Segment.mono}
\difficulty{1}
Let $M$ be a multi-tape machine, $F,F'$ be sets of states, and $P,P'$ be predicates on
work-tape head positions. If a run from $a$ reaches $b$ after $t$ steps, avoids $F$
before time $t$, satisfies $P$ throughout times $0$ through $t$, and $b$'s state avoids
$F$, then, whenever $F'\subseteq F$ and $P(p)\Rightarrow P'(p)$ for every head-position
vector $p$, the same run is a safe segment relative to $F'$ and $P'$.
\end{theorem}

\begin{theorem}[Safe segment endpoint avoids forbidden states]
\lean{Turing.SafeSegment.endpoint_not_mem}
\label{Turing.SafeSegment.endpoint_not_mem}
\leanok
\uses{Turing.Cfg, Turing.SafeSegment}
\difficulty{1}
For a multi-tape machine, a set $F$ of forbidden states, a predicate $P$ on work-tape
head positions, configurations $a,b$, and a natural number $t$, if the run from $a$
reaches $b$ after $t$ steps, avoids $F$ at every time strictly before $t$, satisfies $P$
at every time from $0$ through $t$, and $b$'s state is outside $F$, then $b$'s state is
outside $F$.
\end{theorem}

\begin{theorem}[A Contract Run Forms a Segment]
\lean{Turing.Contract.segment}
\label{Turing.Contract.segment}
\leanok
\uses{Turing.Cfg, Turing.Contract, Turing.HeadBox, Turing.Segment}
\difficulty{1}
Let $C$ be a contract for a machine $M$ with entry state $e$ and exit set $F$. For any
data value $i$ and configuration $d$ in state $e$ satisfying the contract precondition,
the run from $d$ reaches the prescribed endpoint at time $C$’s assigned time, its state
lies outside $F$ at every strictly earlier time, and at every time up to and including
the assigned time each work-tape head position lies between the corresponding lower and
upper bounds specified by $C$.
\end{theorem}

\begin{definition}[Contract from a computation segment]
\lean{Turing.Contract.ofSegment}
\label{Turing.Contract.ofSegment}
\leanok
\uses{Turing.Cfg, Turing.Contract, Turing.HeadBox, Turing.Segment}
Let $M$ be a $k$-work-tape Turing machine with entry state $e$ and a designated set of
exit states, and fix an input and an index type $I$. Suppose a precondition, running
time, endpoint configuration, and lower and upper work-head bounds are specified for
each index and starting configuration. If, whenever the starting configuration is in
state $e$ and satisfies the precondition, the computation segment of the specified
duration reaches the specified endpoint, avoids the exit set at all earlier times, and
keeps every work-tape head between its specified bounds throughout the segment, and the
endpoint state belongs to the exit set, then these data define a contract with those
preconditions, times, endpoints, and bounds.
\end{definition}

\begin{theorem}[Work-Tape Head Position Bound]
\lean{Turing.Contract.time_envelope}
\label{Turing.Contract.time_envelope}
\leanok
\uses{Turing.Cfg, Turing.Contract, Turing.MultiTapeTM.runFrom, Turing.MultiTapeTM.workTapePos_runFrom_mem_Icc}
\difficulty{3}
Let $M$ be a machine and $C$ a contract for $M$. For any data value $i$, configuration
$d$, natural number $u$ with $u\le C$’s assigned running time at $(i,d)$, and work-tape
index $j$, the head position of tape $j$ after $u$ steps lies in the interval from its
initial position minus $C$’s assigned running time to its initial position plus that
running time, inclusive.
\end{theorem}

\begin{definition}[Round stopping states]
\lean{Turing.roundStop}
\label{Turing.roundStop}
\leanok
For a state $a$ and a set $E$ of optional states, the round stopping set consists of all
states in $E$, the state $a$, and the halting state.
\end{definition}

\begin{theorem}[Characterizing states outside the stop set]
\lean{Turing.not_mem_roundStop}
\label{Turing.not_mem_roundStop}
\leanok
\uses{Turing.roundStop}
\difficulty{4}
For a set $E\subseteq\mathrm{Option}(S)$, element $a\in S$, and
$q\in\mathrm{Option}(S)$, $q$ lies outside the set consisting of the absent state, $a$
as a present state, and all elements of $E$ if and only if there is some $s\in S$ such
that $q$ is the present state $s$, $s\ne a$, and the present state $s$ does not belong
to $E$.
\end{theorem}

\begin{definition}[Turing machine round]
\lean{Turing.Round}
\label{Turing.Round}
\leanok
\uses{Turing.Cfg, Turing.MultiTapeTM, Turing.MultiTapeTM.runFrom, Turing.roundStop}
Let $M$ be a $k$-work-tape Turing machine, let $a$ be a state, let $E$ be a set of
optional states, and let $P$ be a predicate on work-tape head positions. For
configurations $c$ and $c'$ on a fixed input and a natural number $t$, require that $c$
starts in state $a$, that $t>0$, and that after exactly $t$ steps the machine reaches
$c'$ with state in $\bigl\{a\bigr\}\cup E$. At every intermediate time $u$ with $0<u<t$,
the state must lie outside the union of $E$, the state $a$, and the halting state;
throughout times $0$ through $t$, the work-tape head positions must satisfy $P$.
\end{definition}

\begin{theorem}[A round is a segment]
\lean{Turing.Round.segment}
\label{Turing.Round.segment}
\leanok
\uses{Turing.Cfg, Turing.MultiTapeTM.runFrom_zero, Turing.Round, Turing.Segment}
\difficulty{4}
Let $M$ be a $k$-work-tape Turing machine, $a$ a state, $E$ a set of optional states,
and $P$ a property of work-head positions. If $c$ and $c'$ form a round of duration $t$
and $\mathrm{some}(a)\notin E$, then the run from $c$ to $c'$ is a segment of duration
$t$: before the endpoint its state avoids $E$, and through the endpoint its work-head
positions satisfy $P$.
\end{theorem}

\begin{theorem}[Intermediate states in a round]
\lean{Turing.Round.cut_some}
\label{Turing.Round.cut_some}
\leanok
\uses{Turing.Cfg, Turing.MultiTapeTM.runFrom, Turing.Round, Turing.not_mem_roundStop}
\difficulty{1}
Let $M$ be a $k$-work-tape Turing machine, let $a$ be a state, let $E$ be a set of
optional states, and let $P$ be a property of work-head positions. If $c$ and $c'$ form
a round of duration $t$, then at every time $u$ with $0<u<t$, the machine is in some
state $s$ with $s
e a$ and whose optional state is not in $E$.
\end{theorem}

\begin{theorem}[A round from an intermediate-state cut]
\lean{Turing.Round.of_cut}
\label{Turing.Round.of_cut}
\leanok
\uses{Turing.Cfg, Turing.MultiTapeTM.runFrom, Turing.Round, Turing.not_mem_roundStop}
\difficulty{4}
Let $M$ be a $k$-work-tape Turing machine, let $a,e$ be states, and let $c$ be a
configuration. Suppose $c$ is in state $a$, $t>0$, and at each time $u$ with $0<u<t$ the
machine is in a state $s$ different from both $a$ and $e$. If at time $t$ the machine is
in state $a$ or $e$, then the run from $c$ to its configuration at time $t$ is a round
from $a$ with escape set consisting of the optional state $e$ and with the head-position
property that holds everywhere.
\end{theorem}

\begin{theorem}[Head positions remain within a time box]
\lean{Turing.Round.time_box}
\label{Turing.Round.time_box}
\leanok
\uses{Turing.Cfg, Turing.HeadBox, Turing.MultiTapeTM.workTapePos_runFrom_mem_Icc, Turing.Round}
\difficulty{4}
Let $M$ be a $k$-work-tape Turing machine, $a$ a state, $E$ a set of optional states,
and $P$ a property of work-head positions. If $c$ and $c'$ form a round of duration $t$,
then they also form a round of duration $t$ in which, at every time $u\le t$, each
work-tape head position lies between its initial position minus $t$ and its initial
position plus $t$, inclusive.
\end{theorem}

\begin{definition}[Orbit under repeated computation rounds]
\lean{Turing.roundOrbit}
\label{Turing.roundOrbit}
\leanok
\uses{Turing.Cfg, Turing.MultiTapeTM, Turing.MultiTapeTM.runFrom}
Let $M$ be a $k$-work-tape Turing machine on a fixed input, and assign a natural-number
duration to each configuration. Starting from a configuration $c$, the configuration
after $r$ rounds is obtained by repeating $r$ times the operation that runs $M$ from the
current configuration for its assigned duration.
\end{definition}

\begin{theorem}[Successive rounds of an orbit]
\lean{Turing.roundOrbit_succ}
\label{Turing.roundOrbit_succ}
\leanok
\uses{Turing.Cfg, Turing.MultiTapeTM, Turing.MultiTapeTM.runFrom, Turing.roundOrbit}
\difficulty{1}
Let $k$ be a natural number, let $\Gamma$ and $S$ be types, let $M$ be a $k$-tape
machine over $\Gamma$ with states in $S$, and let $\tau$ assign a natural number of
steps to each configuration on an input word $x$. For any such configuration $c$ and
natural number $r$, the configuration after $r+1$ rounds is obtained by running $M$ for
$\tau$ of the configuration after $r$ rounds, starting from that configuration.
\end{theorem}

\begin{theorem}[Concatenating successive rounds]
\lean{Turing.Round.iterate}
\label{Turing.Round.iterate}
\leanok
\uses{Turing.Cfg, Turing.MultiTapeTM, Turing.Round, Turing.Round.segment, Turing.Segment, Turing.Segment.trans, Turing.Segment.zero, Turing.roundOrbit}
\difficulty{4}
Let $M$ be a multi-tape Turing machine, $a$ a state, $E$ a set of optional states, and
$P$ a property of work-tape head positions. Let $\tau$ assign a natural-number duration
to each configuration, and let $c$ be a configuration. Assume that the optional state
corresponding to $a$ is not in $E$, and that $P$ holds at the work-tape head positions
of $c$. Define $c_r$ by starting at $c$ and, $r$ times, running the machine for the
duration assigned by $\tau$ to the current configuration. For every $r<n$, assume that
$c_r$ is in state $a$, that its assigned duration is positive, and that running it for
that duration reaches $c_{r+1}$, whose state is either $a$ or belongs to $E$; at each
strictly intermediate positive time the state is neither halted nor in $E$ nor equal to
$a$, and $P$ holds at every time from the start through the endpoint. Then running the
machine from $c$ for the sum of these $n$ durations reaches $c_n$; at every time
strictly before the endpoint its state is not in $E$, and $P$ holds at every time from
the start through the endpoint.
\end{theorem}

\begin{definition}[Round contract]
\lean{Turing.RoundContract}
\label{Turing.RoundContract}
\leanok
\uses{Turing.Cfg, Turing.MultiTapeTM, Turing.MultiTapeTM.runFrom, Turing.roundStop}
Let $M$ be a $k$-work-tape Turing machine with state set $S$, anchor state $a$, escape
set $E\subseteq \{\mathrm{None}\}\cup\{\mathrm{Some}(s):s\in S\}$, index type $I$, and
fixed input $x$. A round contract consists of a precondition, a time, an endpoint
configuration, and lower and upper integer bounds for each work-tape head, all depending
on an index and a starting configuration. Whenever a starting configuration is in state
$a$ and satisfies the precondition, its time is positive, the machine reaches the
specified endpoint after exactly that many steps, and the endpoint state is either $a$
or an element of $E$. At every positive time strictly before the endpoint, the machine's
state is neither halted nor in the anchor-or-escape stopping set; at every time up to
and including the endpoint, each work-tape head position lies between its specified
lower and upper bounds, inclusive.
\end{definition}

\begin{theorem}[A Round Satisfies Its Contract]
\lean{Turing.RoundContract.round}
\label{Turing.RoundContract.round}
\leanok
\uses{Turing.Cfg, Turing.HeadBox, Turing.Round, Turing.RoundContract}
\difficulty{1}
Let $M$ be a $k$-tape Turing machine with state set $S$, let $a\in S$, let $E$ be a set
of optional states, and let $R$ be a round contract for $M$ with anchor $a$ and escape
set $E$. For every data value $i$ and configuration $d$ whose state is $a$ and which
satisfies the contract's precondition, the run for the prescribed duration is a round:
the duration is positive, its endpoint is the prescribed post-configuration with state
in the anchor state or $E$, no strictly intermediate positive time has a halted, anchor,
or escape state, and all work-tape head positions throughout the run lie within the
contract's coordinatewise bounds.
\end{theorem}

\begin{definition}[Returning contract from a round contract]
\lean{Turing.RoundContract.returning}
\label{Turing.RoundContract.returning}
\leanok
\uses{Turing.ReturnContract, Turing.RoundContract}
Let $R$ be a round contract for a machine with anchor state $a$ and escape set $E$.
Restrict its precondition to starting configurations whose specified endpoint has state
$a$. With the same running times, endpoints, and head bounds, this gives a return
contract: for every such start in state $a$, the duration is positive, the computation
reaches the endpoint at that time, the endpoint is in state $a$, the computation does
not revisit $a$ at any positive time strictly before then, and the head bounds hold
throughout.
\end{definition}

\begin{definition}[Escaping contract from a round contract]
\lean{Turing.RoundContract.escaping}
\label{Turing.RoundContract.escaping}
\leanok
\uses{Turing.Contract, Turing.Round.segment, Turing.RoundContract, Turing.RoundContract.round}
Let $R$ be a round contract for a machine with anchor state $a$ and escape set $E$, and
suppose the state $a$ is not in $E$. Restrict the precondition to starting
configurations whose specified endpoint state belongs to $E$. With the same running
times, endpoints, and head bounds, this gives a contract with entry state $a$ and exit
set $E$: from every such start in state $a$, the computation reaches the endpoint at the
specified time, its state belongs to $E$, the computation avoids $E$ at all earlier
times, and the head bounds hold throughout.
\end{definition}

\begin{definition}[Return contract as a round contract]
\lean{Turing.ReturnContract.toRound}
\label{Turing.ReturnContract.toRound}
\leanok
\uses{Turing.MultiTapeTM, Turing.MultiTapeTM.runFrom, Turing.MultiTapeTM.runFrom_state_ne_none_of_le, Turing.ReturnContract, Turing.RoundContract, Turing.not_mem_roundStop}
Given a $k$-tape machine, an anchor state, an index type, an input word, and a return
contract for that machine and anchor, the associated round contract has the same
precondition, running time, endpoint, and work-head bounds. For every eligible start
configuration, its run takes positive time, returns to the anchor, avoids both the
anchor and halting state at every strictly intermediate positive time, and keeps every
work-head position within the specified bounds throughout the run.
\end{definition}

\begin{definition}[Lower Endpoint of a Round-Orbit Head Range]
\lean{Turing.RoundContract.hullLo}
\label{Turing.RoundContract.hullLo}
\leanok
\uses{Turing.Cfg, Turing.RoundContract, Turing.roundOrbit}
Let $R$ be a round contract for a $k$-tape Turing machine $M$, let $i$ be an index in
its index type, let $c$ be a configuration, let $n$ be a nonnegative integer, and let
$j$ be a work-tape index. Write $c_r$ for the configuration obtained by iterating the
round orbit from $c$ exactly $r$ times, using the duration specified by $R$. The lower
endpoint is the minimum of the prescribed lower bounds $R.lo(i,c_r,j)$ for $0 \le r < n$
together with the position of the head on tape $j$ in $c_n$.
\end{definition}

\begin{definition}[Upper endpoint of a round-orbit envelope]
\lean{Turing.RoundContract.hullHi}
\label{Turing.RoundContract.hullHi}
\leanok
\uses{Turing.Cfg, Turing.RoundContract, Turing.roundOrbit}
Let $R$ be a round contract for a $k$-tape machine, let $i$ be an index, let $c$ be a
configuration, and let $n$ be a nonnegative integer. For a work-tape index $j$, define
$H_j$ to be the maximum of the upper envelope values assigned by $R$ to the
configurations after $r$ rounds from $c$, for $0 \le r < n$, together with the head
position on tape $j$ after $n$ rounds.
\end{definition}

\begin{theorem}[Anchor States Along a Round Orbit]
\lean{Turing.RoundContract.orbit_anchor}
\label{Turing.RoundContract.orbit_anchor}
\leanok
\uses{Turing.Cfg, Turing.RoundContract, Turing.roundOrbit, Turing.roundOrbit_succ}
\difficulty{4}
Let $M$ be a multi-tape Turing machine with state set $S$, let $a\in S$, and let $R$ be
a round contract for $M$ at anchor $a$. If $c$ starts in state $a$, the precondition
holds at every orbit configuration with index $r<n$, and each round endpoint with
$r+1<n$ returns to state $a$, then every orbit configuration with index $r<n$ is in
state $a$.
\end{theorem}

\begin{theorem}[Rounds Along a Round Orbit]
\lean{Turing.RoundContract.orbit_round}
\label{Turing.RoundContract.orbit_round}
\leanok
\uses{Turing.Cfg, Turing.HeadBox, Turing.Round, Turing.RoundContract, Turing.RoundContract.hullHi, Turing.RoundContract.hullLo, Turing.RoundContract.orbit_anchor, Turing.RoundContract.round, Turing.roundOrbit, Turing.roundOrbit_succ}
\difficulty{6}
Let $M$ be a $k$-tape Turing machine with state set $S$, let $a\in S$, let $E$ be a set
of optional states, and let $R$ be a round contract for $M$ with anchor $a$ and escape
set $E$. Suppose $c$ starts in state $a$, the precondition holds at every orbit
configuration with index $r<n$, and each round endpoint with $r+1<n$ returns to state
$a$. Then for every $r<n$, the run from orbit configuration $r$ to orbit configuration
$r+1$ is a round of duration prescribed by $R$, with endpoint in the anchor state or
$E$, no halting, anchor, or escape state at strictly intermediate positive times, and
all work-tape head positions within the coordinatewise bounds given by the minimum and
maximum of the contract bounds at orbit indices below $n$ together with the head
positions at orbit index $n$.
\end{theorem}

\begin{theorem}[Initial Heads in the Orbit Envelope]
\lean{Turing.RoundContract.orbit_start}
\label{Turing.RoundContract.orbit_start}
\leanok
\uses{Turing.Cfg, Turing.HeadBox, Turing.RoundContract, Turing.RoundContract.hullHi, Turing.RoundContract.hullLo, Turing.RoundContract.orbit_round, Turing.roundOrbit}
\difficulty{4}
Let $M$ be a $k$-tape Turing machine, let $a\in S$, and let $R$ be a round contract for
$M$ with anchor $a$. Suppose $c$ starts in state $a$, the precondition holds at every
orbit configuration with index $r<n$, and each round endpoint with $r+1<n$ returns to
state $a$. Then the vector of work-tape head positions in $c$ lies in the coordinatewise
box whose lower and upper endpoints are, respectively, the minimum and maximum of the
contract bounds at orbit indices below $n$ together with the head positions at orbit
index $n$.
\end{theorem}

\begin{theorem}[Head-box bounds on the last coordinate]
\lean{Turing.headBox_lastCases}
\label{Turing.headBox_lastCases}
\leanok
\uses{Turing.HeadBox}
\difficulty{4}
Let $k$ be a natural number, let $lo,hi:\mathrm{Fin}(k)\to\mathbb{Z}$, let
$clo,chi\in\mathbb{Z}$, and let $p:\mathrm{Fin}(k+1)\to\mathbb{Z}$. Suppose the last
coordinate of $p$ lies in $[clo,chi]$, and each of its first $k$ coordinates lies in the
corresponding interval $[lo_j,hi_j]$. Then every coordinate of $p$ lies in the
corresponding interval whose bounds are formed by adjoining $clo$ and $chi$ as the
last-coordinate bounds.
\end{theorem}

\begin{definition}[Predicate Along a Computation]
\lean{Turing.Contract.Traj}
\label{Turing.Contract.Traj}
\leanok
\uses{Turing.Cfg, Turing.Contract, Turing.MultiTapeTM, Turing.MultiTapeTM.runFrom}
Let $C$ be a contract for a multi-tape Turing machine, with index set $\iota$, fixed
input, entry state $e$, and prescribed running time. Let $R(i,d,p)$ be a predicate on an
index $i$, a starting configuration $d$, and a vector $p$ of work-tape head positions.
For every $i\in\iota$ and starting configuration $d$ whose state is $e$ and which
satisfies the contract precondition, $R$ holds of the vector of head positions after $u$
steps for every natural number $u$ with $u\leq C$’s prescribed running time on $i,d$.
\end{definition}

\begin{theorem}[Trajectory Within the Head Envelope]
\lean{Turing.Contract.traj_env}
\label{Turing.Contract.traj_env}
\leanok
\uses{Turing.Contract, Turing.Contract.Traj, Turing.HeadBox}
\difficulty{1}
Let $M$ be a multi-tape Turing machine and $C$ a contract for $M$, with lower and upper
head bounds for each work tape. For every parameter $i$ and starting configuration $d$
in the contract entry state and satisfying its precondition, at every time $u$ up to and
including the prescribed running time, the vector of work-tape head positions lies in
the box determined by those bounds; that is, each head position lies between its
corresponding lower and upper bounds.
\end{theorem}

\begin{theorem}[Monotonicity of trajectory properties]
\lean{Turing.Contract.Traj.mono}
\label{Turing.Contract.Traj.mono}
\leanok
\uses{Turing.Cfg, Turing.Contract, Turing.Contract.Traj}
\difficulty{1}
Let $C$ be a contract, and let $R,R'$ be predicates on its data, starting
configurations, and work-head position vectors. If $R$ holds throughout every eligible
trajectory of $C$ and $R(i,d,p)$ implies $R'(i,d,p)$ for every datum $i$, configuration
$d$, and head-position vector $p$, then $R'$ also holds throughout every eligible
trajectory of $C$.
\end{theorem}

\begin{theorem}[Conjunction of trajectory properties]
\lean{Turing.Contract.Traj.and}
\label{Turing.Contract.Traj.and}
\leanok
\uses{Turing.Cfg, Turing.Contract, Turing.Contract.Traj}
\difficulty{1}
Let $C$ be a contract, and let $R$ and $R'$ be predicates on its data, starting
configurations, and work-head position vectors. If both predicates hold at every time up
to and including the contract duration for every starting configuration in the entry
state satisfying the precondition, then their conjunction also holds at every such time.
\end{theorem}

\begin{theorem}[Trajectory property gives a computation segment]
\lean{Turing.Contract.Traj.segment}
\label{Turing.Contract.Traj.segment}
\leanok
\uses{Turing.Cfg, Turing.Contract, Turing.Contract.Traj, Turing.Segment}
\difficulty{1}
Let $C$ be a contract for a machine $M$ with exit set $F$, and suppose $R$ holds
throughout its trajectories. For any datum $i$ and starting configuration $d$ in the
entry state satisfying the precondition, the run from $d$ for the prescribed duration
reaches the prescribed endpoint, its state avoids $F$ at every strictly earlier time,
and $R(i,d,p)$ holds for the work-head position vector $p$ at every time up to and
including that duration.
\end{theorem}

\begin{theorem}[Trajectory under contract composition]
\lean{Turing.Contract.Traj.trans}
\label{Turing.Contract.Traj.trans}
\leanok
\uses{Turing.Cfg, Turing.Contract, Turing.Contract.Traj, Turing.Contract.trans, Turing.MultiTapeTM.runFrom, Turing.MultiTapeTM.runFrom_add}
\difficulty{4}
Let $M$ be a machine, let $C_1$ be a contract from entry state $e_1$ to the sole exit
state $e_2$, and let $C_2$ be a contract from $e_2$ to an exit set $X$. Suppose a map
sends each datum $i$ for $C_1$ to a datum for $C_2$, every admissible endpoint of $C_1$
satisfies the corresponding precondition of $C_2$, and the run of $M$ from any
admissible $C_1$ start has state outside $X$ at every time strictly before $C_1$'s
assigned time. If $R_1$ holds along every admissible $C_1$ trajectory and $R_2$ holds
along every admissible $C_2$ trajectory, then at every time up to and including the
composed contract's assigned time, the work-tape head-position vector satisfies either
$R_1(i,d,p)$ or $R_2$ for the linked datum and the endpoint of $C_1$, at that same
vector $p$.
\end{theorem}

\begin{theorem}[Frames under contract composition]
\lean{Turing.Contract.Frames.trans}
\label{Turing.Contract.Frames.trans}
\leanok
\uses{Turing.Cfg, Turing.Cfg.SameOutside.trans, Turing.Contract, Turing.Contract.Frames, Turing.Contract.trans, Turing.MultiTapeTM.runFrom}
\difficulty{3}
Let $M$ be a machine and let $C_1$ be a contract from entry state $e_1$ to the sole exit
state $e_2$, and $C_2$ a contract from $e_2$ to an exit set $X$, on the same input with
the same work tapes. Suppose a map from the data for $C_1$ to the data for $C_2$ makes
every eligible endpoint of $C_1$ satisfy the corresponding precondition of $C_2$, and
the run under $C_1$ avoids $X$ strictly before its prescribed endpoint. If the two
contracts preserve tape contents outside their respective frame sets, then the composed
contract preserves tape contents outside the union of the first frame set and the second
frame set evaluated at the linked datum and the first endpoint.
\end{theorem}

\begin{theorem}[Zero Running Time and Exit Membership]
\lean{Turing.Contract.time_eq_zero_iff}
\label{Turing.Contract.time_eq_zero_iff}
\leanok
\uses{Turing.Cfg, Turing.Contract, Turing.MultiTapeTM.runFrom_zero}
\difficulty{4}
Let $M$ be a multi-tape Turing machine, let $C$ be a contract for $M$ with entry state
$e$ and exit set $E$, and let $i$ and $d$ be a parameter and starting configuration. If
$d$ is in state $e$ and satisfies the precondition of $C$ for $i$, then the prescribed
running time is zero if and only if the state $e$ belongs to $E$.
\end{theorem}

\begin{theorem}[Excluding Entry from Exits Gives Positive Duration]
\lean{Turing.Contract.positive_of_not_mem}
\label{Turing.Contract.positive_of_not_mem}
\leanok
\uses{Turing.Contract, Turing.Contract.Positive, Turing.Contract.time_eq_zero_iff}
\difficulty{1}
Let $C$ be a contract with entry state $e$ and exit set $X$. If
$\mathrm{some}(e)\notin X$, then every admissible start configuration in state $e$ has
strictly positive assigned running time.
\end{theorem}

\begin{theorem}[Positive Duration Excludes the Entry State]
\lean{Turing.Contract.not_mem_of_positive}
\label{Turing.Contract.not_mem_of_positive}
\leanok
\uses{Turing.Cfg, Turing.Contract, Turing.Contract.Positive, Turing.Contract.time_eq_zero_iff}
\difficulty{1}
Let $C$ be a contract with entry state $e$ and exit set $X$. If every admissible start
configuration in state $e$ has strictly positive assigned running time, then the
optional state $\mathrm{some}(e)$ does not belong to $X$.
\end{theorem}

\begin{theorem}[Positivity of the Composed Running Time]
\lean{Turing.Contract.trans_positive}
\label{Turing.Contract.trans_positive}
\leanok
\uses{Turing.Contract, Turing.Contract.Positive, Turing.Contract.trans, Turing.MultiTapeTM.runFrom}
\difficulty{1}
Let $C_1$ be a contract for $M$ from entry state $e_1$ to the singleton exit state
$e_2$, and let $C_2$ be a contract for $M$ from $e_2$ to exit set $X$. Let
$\mathrm{link}$ map parameters of $C_1$ to parameters of $C_2$. Suppose every endpoint
of $C_1$ from an admissible start satisfies the corresponding linked precondition of
$C_2$, and before the end of the first run the state is outside $X$. If $C_1$ has
strictly positive prescribed running time for every admissible start, then the contract
obtained by composing $C_1$ and $C_2$ has strictly positive prescribed running time for
every admissible start.
\end{theorem}

\begin{lemma}[Trajectory under State Embedding]
\lean{Turing.Contract.mapState_traj}
\label{Turing.Contract.mapState_traj}
\leanok
\uses{Turing.Action.mapState, Turing.Cfg, Turing.Cfg.mapState, Turing.Contract, Turing.MultiTapeTM, Turing.MultiTapeTM.runFrom, Turing.MultiTapeTM.runFrom_mapState_of_agreeOn}
\difficulty{4}
Let $M$ and $N$ be machines with the same number of work tapes and input alphabet, with
state sets $S$ and $H$, and let $C$ be a contract for $M$ with entry state $e$ and exit
set $X$. Let $\mathrm{emb}:S\hookrightarrow H$ be an injective map, and let
$\mathrm{good}$ be a property of states in $S$. Suppose that whenever $q$ is good, the
transition of $N$ at $\mathrm{emb}(q)$ is the transition of $M$ at $q$ with its
successor state mapped by $\mathrm{emb}$. Suppose also that, for every data value $i$
and admissible start configuration $c$ in state $e$, every state visited by $M$ at a
time strictly less than $C$’s assigned time is good. Then for every $i$ and every
configuration $d$ of $N$ whose state is $\mathrm{emb}(e)$, if the configuration
obtained from $d$ by replacing its state with $e$ satisfies $C$’s precondition, the run
of $N$ from $d$ agrees, at every time up to and including that assigned time, with the
run of $M$ from the state-replaced configuration after mapping its state by
$\mathrm{emb}$.
\end{lemma}

\begin{definition}[Transporting a contract through a state embedding]
\lean{Turing.Contract.mapState}
\label{Turing.Contract.mapState}
\leanok
\uses{Turing.Action.mapState, Turing.Cfg.mapState, Turing.Contract, Turing.Contract.mapState_traj, Turing.MultiTapeTM, Turing.MultiTapeTM.runFrom}
Let $C$ be a contract for a machine $M$ with entry state $e$ and exit set $X$, let $N$
be another machine, and let $f$ be an injective map from the states of $M$ to those of
$N$. Suppose that, at every state satisfying a predicate $G$, the transition of $N$ at
its image under $f$ is the state-renamed transition of $M$; also suppose that every
state visited by $M$ at a time strictly before $C$’s endpoint satisfies $G$, for every
eligible start. Then there is a contract for $N$ with entry state $f(e)$ and exit set
consisting of the images under $f$ of the states in $X$: its start configurations are
decoded by replacing their optional state with $e$, and its endpoint states are mapped
by $f$, while its running times and work-head bounds are inherited from $C$.
\end{definition}

\begin{definition}[Lockstep simulation between machines]
\lean{Turing.Lockstep}
\label{Turing.Lockstep}
\leanok
\uses{Turing.Cfg, Turing.MultiTapeTM, Turing.MultiTapeTM.runFrom, Turing.MultiTapeTM.step}
Let $M$ and $N$ be multi-tape Turing machines, let $e$ map configurations of $M$ to
configurations of $N$, and let $c\mapsto t(c)$ assign a nonnegative integer to each
configuration $c$ of $M$. Let $X$ and $X'$ be sets of states, allowing the halting state
to be represented by absence of a state. If the state of $c$ is outside $X$, then
running $N$ from $e(c)$ for $t(c)$ steps yields the encoding of the configuration
reached by one step of $M$ from $c$; at every time $v<t(c)$, the state reached by $N$ is
outside $X'$. If the state of $c$ belongs to $X$, then the state of $e(c)$ belongs to
$X'$.
\end{definition}

\begin{theorem}[Simulation of a Run by Total Cost]
\lean{Turing.Lockstep.runFrom}
\label{Turing.Lockstep.runFrom}
\leanok
\uses{Turing.Cfg, Turing.Lockstep, Turing.MultiTapeTM.runFrom, Turing.MultiTapeTM.runFrom_add, Turing.MultiTapeTM.runFrom_succ_eq_step'}
\difficulty{4}
Let $M$ and $N$ be multi-tape Turing machines, let $e$ map configurations of $M$ to
configurations of $N$, let $q$ assign a natural-number cost to each configuration of
$M$, and let $X$ and $X'$ be sets of optional states of $M$ and $N$, respectively.
Suppose that whenever a configuration $c$ of $M$ has state outside $X$, running $N$ from
$e(c)$ for $q(c)$ steps yields $e$ of the one-step successor of $c$ under $M$;
configurations with state in $X$ encode to configurations whose state lies in $X'$, and
every run of $N$ from $e(c)$ stays outside $X'$ at all times strictly less than $q(c)$
when $c$ has state outside $X$. If a configuration $d$ of $M$ has state outside $X$ at
each time $u<t$ in its run, then running $N$ from $e(d)$ for the sum of the costs of the
first $t$ configurations in that run yields the encoding of the configuration reached by
$M$ after $t$ steps.
\end{theorem}

\begin{theorem}[Avoidance before accumulated simulation cost]
\lean{Turing.Lockstep.cut}
\label{Turing.Lockstep.cut}
\leanok
\uses{Turing.Cfg, Turing.Lockstep, Turing.Lockstep.runFrom, Turing.MultiTapeTM.runFrom, Turing.MultiTapeTM.runFrom_add}
\difficulty{5}
Let $M$ and $N$ be multi-tape Turing machines, let $e$ map configurations of $M$ to
configurations of $N$, let $q$ assign a natural-number simulation cost to each
configuration of $M$, and let $X$ and $X'$ be sets of states (including the possible
halting state) of $M$ and $N$, respectively. Assume that for every configuration $c$
whose state is outside $X$, running $N$ from $e(c)$ for $q(c)$ steps yields the encoding
of the next configuration of $M$; that if $c$ has state in $X$, then $e(c)$ has state in
$X'$; and that, for every such $c$, the run of $N$ from $e(c)$ has state outside $X'$ at
every time strictly less than $q(c)$. If the run of $M$ from an initial configuration
$d$ has state outside $X$ at every time $u<t$, then the run of $N$ from $e(d)$ has state
outside $X'$ at every time strictly less than $\sum_{u=0}^{t-1} q(M^u(d))$, where
$M^u(d)$ denotes the configuration reached after $u$ steps.
\end{theorem}

\begin{theorem}[Locating a simulated configuration]
\lean{Turing.Lockstep.locate}
\label{Turing.Lockstep.locate}
\leanok
\uses{Turing.Cfg, Turing.Lockstep, Turing.Lockstep.runFrom, Turing.MultiTapeTM.runFrom, Turing.MultiTapeTM.runFrom_add}
\difficulty{5}
Let $M$ and $N$ be multi-tape Turing machines, let $e$ map configurations of $M$ to
configurations of $N$, let $q$ assign a natural-number cost to each configuration of
$M$, and let $X$ and $X'$ be sets of states, including the possible halting state, for
$M$ and $N$, respectively. Assume that whenever a configuration $c$ of $M$ has state
outside $X$, running $N$ from $e(c)$ for $q(c)$ steps yields the encoding of the next
configuration of $M$; that states in $X$ encode into states in $X'$; and that each such
run of $N$ stays outside $X'$ at every time strictly less than $q(c)$. If the run of $M$
from a configuration $d$ has state outside $X$ at every time $u<t$, then for every
natural number $v$ at most the sum of the costs of the configurations visited by $M$ at
times $0,\ldots,t-1$, there are natural numbers $u\le t$ and $w$ such that either $w=0$,
or the configuration reached by $M$ at time $u$ has state outside $X$ and $w\le q(c_u)$,
where $c_u$ is that configuration, and the configuration reached by $N$ from $e(d)$
after $v$ steps equals the configuration reached by $N$ from $e(c_u)$ after $w$ steps.
\end{theorem}

\begin{definition}[Contract induced by lockstep simulation]
\lean{Turing.Contract.lockstep}
\label{Turing.Contract.lockstep}
\leanok
\uses{Turing.Cfg, Turing.Contract, Turing.Lockstep, Turing.Lockstep.cut, Turing.Lockstep.locate, Turing.Lockstep.runFrom, Turing.MultiTapeTM, Turing.MultiTapeTM.runFrom}
Let $C$ be a contract for a machine $M$ with entry $e$ and exits $X$, and let a lockstep
simulation encode configurations of $M$ into configurations of a machine $N$, with
target exit set $X'$ and per-configuration cost. Fix a target entry state $e'$ and
integer-valued lower and upper bounds $b_{lo}(c,j)$ and $b_{hi}(c,j)$ for each target
work tape $j$. Suppose that the target head position at every encoded configuration
after $v$ steps lies within these bounds whenever $v=0$, or whenever the source state is
outside $X$ and $v$ is at most its cost. Then there is a contract for $N$, indexed by
pairs $(i,c)$, whose precondition requires a configuration equal to $enc(c)$ with $c$ in
state $e$ and satisfying $C$’s precondition. Its running time is the sum of the
simulation costs over the source run up to (but excluding) $C$’s endpoint; its endpoint
is the encoding of $C$’s endpoint; and its work-head bounds are the minimum and maximum
of the supplied bounds over the source run, including the endpoint.
\end{definition}

\begin{theorem}[Admissibility of a lockstep contract]
\lean{Turing.Contract.lockstep_admissible}
\label{Turing.Contract.lockstep_admissible}
\leanok
\uses{Turing.Cfg, Turing.Contract, Turing.Contract.lockstep, Turing.Lockstep, Turing.MultiTapeTM, Turing.MultiTapeTM.runFrom}
\difficulty{1}
Let $M$ and $N$ be machines, let $E$ map configurations of $M$ to configurations of $N$,
and let $c	ext{ost}$ assign a natural-number cost to each configuration of $M$. Let $C$
be a contract for $M$ with entry state $e$ and exit set $X$, and let $L$ be a lockstep
relation between $M$ and $N$ using $E$, $c\text{ost}$, and exit sets $X,X'$. Let $e'$ be
a state of $N$, and let $b_-,b_+$ assign integer bounds to each configuration of $M$ and
tape of $N$. Assume that for every configuration $c$ of $M$, every time $v$ that is zero
or satisfies $c$'s state is outside $X$ and $v\le c\text{ost}(c)$, and every tape of
$N$, the head position after $v$ steps from $E(c)$ lies in $[b_-(c),b_+(c)]$. If $c$ is
in state $e$ and satisfies $C$'s precondition for parameter $i$, then $E(c)$ is in state
$e'$, and it satisfies the precondition of the lockstep contract built from $C$ and $L$
at parameter $(i,c)$.
\end{theorem}

\begin{theorem}[Admissibility under state mapping]
\lean{Turing.Contract.lockstep_admissible_map}
\label{Turing.Contract.lockstep_admissible_map}
\leanok
\uses{Turing.Cfg, Turing.Contract, Turing.Contract.lockstep, Turing.Contract.lockstep_admissible, Turing.Lockstep, Turing.MultiTapeTM, Turing.MultiTapeTM.runFrom}
\difficulty{3}
Let $M,N$ be machines, let $E$ map configurations of $M$ to configurations of $N$, and
let $c\text{ost}$ assign a natural-number cost to each configuration of $M$. Let $C$ be
a contract for $M$ with entry state $e$ and exit set $X$, and let $L$ be a lockstep
relation between $M$ and $N$ using $E$, $c\text{ost}$, and exit sets $X,X'$. Let $\rho$
map states of $M$ to states of $N$, and suppose the state of $E(c)$ is obtained by
applying $\rho$ to the optional state of every configuration $c$. Let $b_-,b_+$ assign
integer bounds to each configuration of $M$ and tape of $N$, and assume that for every
configuration $c$, every time $v$ that is zero or satisfies $c$'s state is outside $X$
and $v\le c\text{ost}(c)$, and every tape of $N$, the head position after $v$ steps from
$E(c)$ lies in $[b_-(c),b_+(c)]$. If $c$ is in state $e$ and satisfies $C$'s
precondition for parameter $i$, then $E(c)$ is in state $\rho(e)$, and it satisfies the
precondition of the lockstep contract built from $C$ and $L$ at parameter $(i,c)$.
\end{theorem}
